% Tier~B: a certified bounded semantics, not a source compiler.
\subsection{A bounded reversible interpreter: scope and obligations}

\begin{remark}[Status of the extension]\label{rem:tB-status}\label{target:tierB}
Definition~\ref{def:tierb} admits a static menu only. The following is
an additional, conditional semantics for a bounded dynamic interpreter.
It specifies what an implementation certificate must establish; it does
not assert that every source program, or the existing executor, has such
a certificate. In particular a supplied sequence of unitary matrices
without the journal and routing obligations below is insufficient.
Every dynamic source interpretation and implementation bound below is
conditional on the supplied certificate and, for circuit costs, the
specified implementation. The static clock schedule and the algebraic
disposal test do not construct a dynamic certificate. General dynamic
source rewind, including the unconditional source-level extension in
Remark~\ref{problem:remaining}, remains open.
\end{remark}

\begin{definition}[Finite interpreter state]\label{def:tB-state}
Fix a finite quantum data register $Q$, an ordinary-work register $A$
encoded in qubits, a program $P$ of length $\ell\geq1$, and integers
$T\geq0$, $N\geq1$, $d\geq0$, and $w\geq1$.
All are supplied independently of the unknown input.
The interpreter register is
\[
 M=A\,C\,J\,S\,R\,Z\,\Gamma.
\]
Here $C$ contains the program position, direction/mode and finite
control; $J$ has $N$ journal cells of $w$ bits and a cursor;
$S$ has $d$ mark cells and a stack pointer; $R$ retains return addresses,
visit identifiers, saved cursors and any other routing information;
$Z$ is implementation workspace; and $\Gamma$ is a clock. All pointer,
identifier and counter widths are explicit. One mark cell records a
mark identifier, an activation identifier, a visit identifier, the
journal cursor, and a continuation address. Each field fits its declared
width. Repeated visits use distinct live identifiers or a proved reversible
reuse rule. Integers are finite binary strings; overflow is never silently
ignored.

A journal entry identifies a specified total unitary, its sign (forward
or adjoint), its operands and recording cells, and the old bits needed
to invert an ordinary update. A record of a gate includes its actual
matrix phase. No quantum state is copied into a journal: ordinary
basis information is copied coherently, and gates are specified by
classical descriptors. Descriptors needed later are retained or recovered
by proved reversible operations. Reading $Q$ means controlled reversible
copying into machine bits. Positions, modes, mark identifiers, cursors and
budgets may consequently be in superposition. The controlled-reversible
reading of primitive~10 has the preserved-sector scope of
Lemma~\ref{lem:read-product}; subsequent gates changing a read register
do not give a controlled description of the entire run.

All wires are reserved initially. There is no observed measurement,
drop, allocation, external output observation or physical stopping during
the $T$ ticks. Output tapes, faults and timeout indicators are retained
registers. A claim of successful source behavior requires that capacities,
valid mark identifiers, blank-cell conditions and the supplied horizon
suffice on the entire allowed input subspace. Completion on other states
is part of the semantics, not a claim of successful source execution.
\end{definition}

\begin{definition}[Local rules and their certificate]\label{def:tB-cert}
Write $B=QACJSR$ for the logical state without clock or clean workspace.
The set of local rules is finite. A certificate names a linear initial
subspace $\mathcal S_{\rm cert}\subseteq
\mathcal H_B\otimes\operatorname{span}\{\ket0_\Gamma\}$;
in an implementation its initial
embedding also appends $\ket0_Z$. The clock is omitted in the autonomous
case. The reachable-subspace condition in item 5 applies to this entire
linear space, not merely to a set of basis inputs. If capacity overflow,
an exhausted journal or stack, or an invalid mark identifier is reachable
within $T$ ticks, the instance has no certificate for successful source
execution on that subspace.
A local rule $e$ consists of an input projector $D_e$ on $B$ and a
partial isometry $V_e$ with
\[
 V_e^\dagger V_e=D_e,\qquad V_eV_e^\dagger=E_e.
\]
Input projectors for different rules are orthogonal, and so are their
output projectors $E_e$. These are \emph{separate obligations}.
Rules are selected by computational-basis mode, position, journal and
mark fields, including coherently computed predicates. Predicates that
are not already machine bits must first be computed reversibly.
Their computation and eventual erasure are themselves included in the
rules. Quantum gates can act within a selected sector.

Each rule supplies explicit coordinate identifications of its input and
output subspaces with the same finite payload space $\mathcal L_e$,
\[
 I_e^-:\mathcal L_e\longrightarrow\mathcal H_B,
 \quad I_e^+:\mathcal L_e\longrightarrow\mathcal H_B,
 \quad V_e=I_e^+F_e(I_e^-)^\dagger,                 \tag{B1}
\]
where $I_e^-$ and $I_e^+$ are isometries onto
$\operatorname{ran}D_e$ and $\operatorname{ran}E_e$, respectively, and $F_e$ is
a specified total unitary. The coordinate identifications are explicit
basis encodings, not unspecified identifications of abstract spaces.
They include every live journal cell and every retained old value.
The forward ordinary action of $F_e$ requires the stated blank cells;
on other codes its chosen total completion still defines the operation.
The snapshot tick of Definition~\ref{def:joint-tick} is one possible
$F_e$; its inverse is the stated $D_t^\dagger F_t^{-1}$, including its
sampling and snapshot registers, not just the inverse data gate.

A certificate verifies the following instruction meanings on a specified
input subspace and all its reachable states through tick $T$.
\begin{enumerate}
\item A forward rule performs its stated completed step and records the
information needed for its inverse, including control changes. A mark
pushes the tuple of Definition~\ref{def:tB-state} reversibly.
\item A keep is the total XOR operation of Definition~\ref{def:keep};
its controlling machine fields are preserved during that operation.
The subsequent change to undo mode retains its old control information.
\item An undo rule selects the last entry above the named mark. If that
entry denotes $F_e$ in (B1), the undo rule has isometric coordinate
encodings $\bar I_e^+,\bar I_e^-$ onto its input and output sectors,
respectively, and is exactly
\[
 \bar V_e=\bar I_e^- F_e^\dagger(\bar I_e^+)^\dagger.          \tag{B2}
\]
The bars permit different direction and continuation labels, subject to
the following matching requirement on the entire retained payload.
Write $\bar D_e=\bar V_e^\dagger\bar V_e$ and
$\bar E_e=\bar V_e\bar V_e^\dagger$. There are specified isometries
$\Lambda_e^+:\operatorname{ran}E_e\to\operatorname{ran}\bar D_e$
and $\Lambda_e^-:\operatorname{ran}\bar E_e\to\operatorname{ran}D_e$,
onto those sectors, such that
\[
 \bar I_e^+=\Lambda_e^+ I_e^+,
 \qquad I_e^-=\Lambda_e^-\bar I_e^-
 \quad\hbox{on all of }\mathcal L_e.                 \tag{B2a}
\]
Every routing or mode operation taking a forward output sector to its
undo input sector realizes $\Lambda_e^+$ on the whole sector; the
corresponding routing after undo realizes $\Lambda_e^-$ on the whole
undo output sector. These are payload-coordinate identities also for
vectors produced by a keep, not only for ordinary forward outputs.
Consecutive recorded entries must admit a common retained payload space
after identity padding. If the specified isometry $J_i$ routes the output of entry $i$ to
the input of entry $i+1$, the specified encodings obey
$J_i I_i^+=I_{i+1}^-$ on that whole space; the reverse routing obeys
$J_i^\dagger I_{i+1}^-=I_i^+$. All these routings retain every live
record and are among the certified operations. Their costs are charged.
The full operator $F_e^\dagger$,
including its journal and control action, is required. A signed adjoint
entry is undone with $F_e$. An inner rewind is represented by its entire
reversible operation, including its keep and routing, or by its sequence
of elementary signed entries; a compound entry is not a unit-cost gate.
\item Undo continues one elementary entry at a time until the recorded
cursor of the coherently named mark is reached. Equality is a reversible
Boolean test. Changing mode and installing a continuation preserve all
information needed to undo those changes. Popped records are retained
until they have been reversibly cleared. Repeated use of a mark and undo
of an earlier rewind obey these same rules; neither operation licenses
re-execution of different, newly selected instructions.
\item For every $0\leq t\leq T$, every vector in
$W^t\mathcal S_{\rm cert}$ lies in the sum of the asserted rule
domains, and each of its sector components satisfies the assertion for
that rule. Here $W$ is the logical tick specified below; in (B3) use
the input projectors $\proj c_\Gamma\otimes D_{c,e}$ for the
clock-$c$ rules, and in the autonomous case use $D_e$.
In an implementation whose single tick $W$ has the supplied clean
implementation $\widetilde W$ of Theorem~\ref{thm:tB-sound}, all
statements are lifted with the fixed zero workspace factor
$\ket0_Z$, which $\widetilde W$ restores at every tick. Thus this condition includes superpositions over
different rule sectors, every post-keep state and every partial undo.
Formula (B2) and the matching identities (B2a) must hold
on that whole domain, not just on states obtained with an empty keep.
All rules have orthogonal ranges as stipulated above. A stopped logical
output may idle only through a reversible rule, retaining any timing
information it requires.
\end{enumerate}
The identifications in (B1)--(B2a), the domain/range assertions, and
reachable-subspace assertions are mathematical proof obligations for an
instance. They are not established merely by naming an instruction.
For an efficient circuit claim, a total clean circuit realizing all
rules and their complement must additionally be supplied.
\end{definition}

\begin{remark}[Why the range obligation matters]\label{rem:tB-range}
For a map specified only on blank records, the allowed input and its
image need not be invariant under keep. The inverse on that image then
has no prescribed action after keep. A total $F_e$ fixes the algebraic
inverse, but a state outside a claimed ordinary sector need not have the
intended cursor, mark or continuation. Extending $F_e$ alone does not
prove the reachable-subspace obligation. Likewise, retaining an extra
copy of a control decision can prevent erasure: reversing the data gates
while leaving that copy is not reversal of the full forward operation.
\end{remark}

\begin{lemma}[Exact extension criterion]\label{lem:tB-extension}
Let $\mathcal H$ be finite dimensional and let $D_1,\ldots,D_k$ be
orthogonal projectors with pairwise orthogonal ranges. For each $i$ let
$V_i$ be an isometry from $D_i\mathcal H$ into $\mathcal H$, extended
by zero on the orthogonal complement. There exists a unitary $U$ on
$\mathcal H$ with $UD_i=V_i$ for every $i$ if and only if the ranges
of the $V_i$ are pairwise orthogonal. The assertion does not impose a
gate-alphabet or efficiency condition on $U$.
\end{lemma}
\begin{proof}
 \step{tB-ext-necessary}{\assume{Such a $U$ exists.}
 \prove{The ranges of the $V_i$ are pairwise orthogonal.}}
 \begin{proof}\pf\ For $x\in D_i\mathcal H$ and
 $y\in D_j\mathcal H$ with $i\ne j$,
 $\langle V_ix,V_jy\rangle=\langle Ux,Uy\rangle
 =\langle x,y\rangle=0$.\end{proof}
 \step{tB-ext-sufficient}{\assume{The ranges are pairwise orthogonal.}
 \prove{Such a $U$ exists.}}
 \begin{proof}\pf\ The direct sum of the $V_i$ is an isometry on
 $\sum_iD_i\mathcal H$. Its input and output dimensions are equal.
 Specify orthonormal bases of their equal-dimensional orthogonal
 complements and map one basis to the other. Combining the two maps
 gives a unitary with the required restrictions.\end{proof}
 \qedstep
 \begin{proof}\pf\ The two implications establish the equivalence.
 \end{proof}
\end{proof}

\begin{definition}[One fixed tick and bounded idle]\label{def:tB-tick}
Set $D=\sum_eD_e$, $E=\sum_eE_e$ and $V=\sum_eV_e$.
A completion is a \emph{specified} unitary isomorphism
$V_\perp:(\id-D)\mathcal H_B\to(\id-E)\mathcal H_B$,
extended by zero elsewhere. Put $U=V+V_\perp$.
In the circuit version this completion must have its supplied clean
implementation; finite-dimensional existence is no efficiency claim.
One may supply such rules and a completion $U_c$ for each clock value
$c$, for example to reserve clock-indexed recording cells and bounded
idle intervals.
In that case $D_{c,e}$ and $V_{c,e}$ denote the input projector and
partial isometry of the clock-$c$ copy of rule $e$, and $U_c$ is the
specified completion formed from them as above.
Choose $L_{\Gamma}=2^b>T$, let $\Gamma$ have $b$ bits initially
zero, and define the single fixed logical tick
\[
 W=\sum_{c=0}^{L_{\Gamma}-1}\ket{c+1\bmod L_{\Gamma}}\bra c_\Gamma\otimes U_c.
                                                               \tag{B3}
\]
Identical $U_c=U$ give the time-independent microcode special case.
A finite autonomous microcode that is already a total unitary can
instead omit the separate clock ($b=0$); a stated horizon still applies.
In (B3) idle may have $U_c=\id$, but the clock still advances bijectively.
The logical run is exactly $W^T$. When a clean implementation of $W$
is supplied, its implementation workspace $Z$ returns to zero each tick.
There is no amplitude-dependent choice of the number of physical ticks.
Logical completion times, mark visits and undo lengths may differ.
The clock has no wrap during the promised interval; outside it the
cyclic continuation is still part of the unitary specification.
\end{definition}

\begin{theorem}[Conditional dynamic soundness and exact undo]
\label{thm:tB-sound}
Assume Definitions~\ref{def:tB-state}--\ref{def:tB-tick}, including all
certificate obligations, specified completions and the full reachable
input subspace $\mathcal S_{\rm cert}$. In particular assume the
matching-encoding identities (B2a), their whole-sector routing
realizations, and the consecutive-entry identifications of
Definition~\ref{def:tB-cert}, including after every keep.
Then $W^T$ is unitary on all retained logical registers.
If the single tick $W$ itself has a supplied clean implementation
$\widetilde W$ on $B\Gamma Z$ in the sense of
Definition~\ref{def:scratch}, namely
$\widetilde W(\psi\otimes\ket0_Z)=W\psi\otimes\ket0_Z$ for every
$\psi$, then $T$ concatenations give a clean implementation with
\[
 (\widetilde W)^T(\psi\otimes\ket0_Z)
       =W^T\psi\otimes\ket0_Z\quad\hbox{for every }\psi.
\]
No clean implementation of $W$ is inferred here from clean
implementations of the separate $U_c$ alone.
On the certified subspace it executes the specified dynamic instructions.
Each undo together with its specified label routing is the exact adjoint
of its corresponding completed forward operation: on all of
$\mathcal L_e$,
\[
 \Lambda_e^-\bar V_e\Lambda_e^+ I_e^+
       = I_e^-F_e^\dagger.                              \tag{B2b}
\]
This includes arbitrary intervening gates, completed inner rewinds and
unequal ticks of arrival at an undo instruction, provided the certificate
obligations hold after those operations. The statement does not assert
cleanliness or a certificate for every source program. General dynamic
source rewind remains open (Open problem~\ref{prob:tB-compiler}).
\end{theorem}
\begin{proof}
 \step{tB-unit-partial}{$V^\dagger V=D$ and $VV^\dagger=E$.}
 \begin{proof}\pf\ For $e\ne f$, orthogonality of output sectors gives
 $V_e^\dagger V_f=0$, and orthogonality of input sectors gives
 $V_eV_f^\dagger=0$. Sum the defining diagonal identities.\end{proof}
 \step{tB-unit-complete}{The completion exists and $U$ is unitary.}
 \begin{proof}\pf\ The preceding step gives equal ranks of $D,E$.
 Their orthogonal complements have equal finite dimension. Choose and
 specify orthonormal bases on those complements and map one to the other.
 Orthogonality eliminates cross terms, so
 $U^\dagger U=D+(\id-D)=\id$ and $UU^\dagger=\id$.
 This argument alone gives no efficient finite-alphabet synthesis.
 Indeed, when the complementary sector is nonzero, exact synthesis of
every prescribed completion over a fixed finite alphabet is
impossible: each finite circuit over that alphabet, with any finite
number of clean workspace wires, is specified by a finite string of
gate names and wire numerals, so these circuits and their exact
logical restrictions form a countable set; whereas, for fixed unit
vectors in the two complementary sectors, mapping the first to
$e^{i\theta}$ times the second and extending by one fixed basis map
elsewhere gives uncountably many distinct unitary completions of the
same $V$ as $\theta$ ranges over $[0,2\pi)$.
 \end{proof}
 \step{tB-unit-clock}{$W$ and $W^T$ are unitary.}
 \begin{proof}\pf\ Distinct summands of (B3) have orthogonal clock
 inputs and outputs. Therefore $W^\dagger W=WW^\dagger=\id$.
 A finite product of this unitary is unitary. For $b=0$ use the preceding
 step directly. If $\widetilde W$ is supplied as hypothesized, induction
 on $t$ gives $(\widetilde W)^t(\psi\ket0_Z)=W^t\psi\ket0_Z$:
 the preceding tick restores zero workspace, so the defining equality
 applies to $W^t\psi$ at the next tick. This also covers $t=0$.
 This product assertion is the elementary part of soundness.\end{proof}
 \step{tB-unit-inverse}{Undo together with its specified label routing
 reverses the entire local completed step.}
 \begin{proof}
  \step{tB-unit-cancel}{Equation (B2b) holds on all of $\mathcal L_e$;
  the routed undo equals $V_e^\dagger$ on $\operatorname{ran}E_e$.}
  \begin{proof}\pf\ For any $v\in\mathcal L_e$, (B2a) and (B2) give
  \[
  \begin{split}
   \Lambda_e^-\bar V_e\Lambda_e^+ I_e^+v
    &=\Lambda_e^-\bar V_e\bar I_e^+v\\
    &=\Lambda_e^-\bar I_e^-F_e^\dagger v
      =I_e^-F_e^\dagger v=V_e^\dagger I_e^+v.
  \end{split}
  \]
  The second equality uses
  $(\bar I_e^+)^\dagger\bar I_e^+=\id$; the last uses (B1).
  Since $I_e^+$ is onto, this proves the sector equality. Taking
  $v=F_eu$ gives $I_e^-u$, the original forward input. The same
  equality holds for every other $v$, including the post-keep vector,
  by the whole-sector routing hypothesis.\end{proof}
  \step{tB-unit-sequence}{For a recorded sequence
  $F_1,\ldots,F_k$ on a common retained payload, with identity padding
  and matching routing identifications included, the inverse payload action is
  $F_1^\dagger\cdots F_k^\dagger$.}
  \begin{proof}\pf\ For $k=0$ the product is identity. For $k>0$,
  write the vector at the last forward output as $I_k^+v$.
  Step~\stepref{tB-unit-cancel} takes it to
  $I_k^-F_k^\dagger v$. If $k>1$, the certified reverse routing
  $J_{k-1}^\dagger$ takes this to $I_{k-1}^+F_k^\dagger v$.
  Induction on the remaining $k-1$ entries gives
  $I_1^-F_1^\dagger\cdots F_k^\dagger v$.
  Forward routing $J_i I_i^+=I_{i+1}^-$ gives the forward payload
  product $F_k\cdots F_1$ by the same induction. Thus the two
  payload products are adjoints. The reachable-subspace obligation
  justifies every sector component used, including after keeps;
  the identities do not assume $v$ has a blank-record origin.
  An inner completed operation $A^\dagger K A$ has adjoint
  $A^\dagger K^\dagger A$ by multiplication of adjoints. The same
  identity applies to its retained control operations.\end{proof}
  \qedstep
  \begin{proof}\pf\ These identities are operator identities on the
  certified sectors. Inserting gates or a keep changes the vector to
  which the inverse is applied, not the inverse matrix. Returning the
  original vector additionally requires undoing all inserted operations;
  $F^\dagger K F$ need not be identity.\end{proof}
 \end{proof}
 \step{tB-unit-dynamic}{Different arrival times require no observation.}
 \begin{proof}\pf\ At any tick decompose the vector by the orthogonal
 rule-domain projectors. Apply the specified linear map to each component
 and sum. Induction over the fixed $T$ ticks uses the reachable-domain
 obligation at each stage. The selected mark and cursor comparisons are
 machine-register operations. No common number of logical undo entries
 is assumed.\end{proof}
 \qedstep
 \begin{proof}\pf\ The preceding steps prove the stated algebraic and
 certified operational assertions. Verification of domains, disjoint
 ranges, (B2) and the whole-sector matching in (B2a), rather than
 multiplication of unitaries, is the
 substantive implementation obligation.\end{proof}
\end{proof}

\begin{remark}[The physical tick is not reversed by advancing its clock]
An undo-mode application of $W$ is not the global operator $W^\dagger$:
the clock in (B3) still advances. The exact-inverse assertion concerns
the completed logical operation and its explicit encodings, including
all of its destroyed-information records. Reversing a physical tick,
clock included, requires $W^\dagger$. Omitting this distinction would
make the exact-inverse claim false even for a classical counter.
\end{remark}

\subsection{The static specialization and the disposal criterion}

\begin{theorem}[Embedding Tier~A, with its machine registers retained]
\label{thm:tB-A}
Let $P_A$ be a finite fixed-space Tier~A term satisfying all applicable
hypotheses of Theorem~\ref{thm:sA-sound}, with
its supplied expansion into logical atoms, keeps and adjoint logical
atoms $A_0,\ldots,A_{n-1}$ in execution order, each a total unitary on
$L_A=QM_A$, with $A_{n-1}\cdots A_0=\mathcal U_{P_A}$.
This is the logical expansion on $L_A$, not its elementary clean circuit
on $L_AZ$. It includes the original occurrence-labelled records and
any specified continuation atoms. Choose $L_{\Gamma}=2^b>T\geq n$ and a fresh
counter $\Gamma$ initially zero, and let $M_B$ be any disjoint auxiliary
register in a fixed unit vector $\ket{m_B}$. Take
$U_c=A_c\otimes\id_{M_B}$ for $0\leq c<n$ and
$U_c=\id_{L_AM_B}$ for $n\leq c<L_{\Gamma}$.
Then this is a static specialization of the tick form (B3), with no
dynamic mark or undo instruction executed, and for every $\psi_{L_A}$,
\[
 W^T(\psi_{L_A}\otimes\ket{m_B}\ket0_\Gamma)
 =\mathcal U_{P_A}\psi_{L_A}\otimes\ket{m_B}\ket T_\Gamma.
                                                               \tag{B4}
\]
The fixed list already specifies every action; no separate mutable mark
stack is used. This theorem supplies no certificate under
Definition~\ref{def:tB-cert}. Any dynamic source interpretation would
additionally require that certificate, whose general construction remains
open (Open problem~\ref{prob:tB-compiler}).
Here a clean implementation concerns only $Z$ as in
Definition~\ref{def:scratch}; logical cleanliness would restore $M_A$
exactly, in the sense of Definition~\ref{def:clean}.
Classical control alone does \emph{not} imply $M_A$ is clean, nor does it
make an arbitrary implementation's extra records independent of $Q$.
Equation (B4) is an existence statement for this implementation.
\end{theorem}
\begin{proof}
 \step{tB-A-word}{The expanded list has product $\mathcal U_{P_A}$.}
 \begin{proof}\pf\ The logical recursion in Definition~\ref{def:tiera}
 and Theorem~\ref{thm:sA-sound} expands identity as an empty list, an
 atom as itself, and a sequence by concatenation. A marked body expands
 as its logical list, its keep, and the reversed adjoint logical list
 using the same record names. Structural induction gives the displayed
 product on all of $\mathcal H_{L_A}$. This never substitutes the
 action of an elementary implementation on nonzero workspace.
 Distinct counter inputs and outputs make the summands of (B3)
 orthogonal. No dynamic certificate is asserted for this schedule.
 \end{proof}
 \step{tB-A-counter}{After $t\leq T$ ticks the clock equals $t$ and
 the logical action is $U_{t-1}\cdots U_0$.}
 \begin{proof}\pf\ Induct on $t$ using (B3), starting with the empty
 product. The assumption $T<L_{\Gamma}$ excludes wrap.\end{proof}
 \qedstep
 \begin{proof}\pf\ At $T$ the last $T-n$ factors are identity.
 Combine the preceding steps. Nothing erases $M_A$; its cleanliness is
 the separate condition in Theorem~\ref{thm:sA-clean}.\end{proof}
\end{proof}

\begin{theorem}[Exact synchronization and disposal]\label{thm:tB-clean}
Let $W$ be a finite unitary interpreter tick, let $T\geq0$, initialize
all of $M$ in a fixed pure state $m_0$, and allow every input on $Q$.
For a dynamic source application, $W$ and its interpretation must be
supplied under the certificate hypotheses of Theorem~\ref{thm:tB-sound}.
To claim that interpretation for every input considered here, under
the identification $QM=B\Gamma Z$ of
Definitions~\ref{def:tB-state}--\ref{def:tB-cert} with factors placed
by name, the fixed vector $m_0$ must have clock factor
$\ket0_\Gamma$ and workspace factor $\ket0_Z$, the tick must be the
logical tick $W$ of Theorem~\ref{thm:tB-sound} (respectively its
supplied clean implementation $\widetilde W$ when $Z$ is retained
in $M$), and $E_0\mathcal H_Q$, with the $Z$ factor suppressed as
in Definition~\ref{def:tB-cert}, must be contained in
$\mathcal S_{\rm cert}$.
The algebraic test below applies to any such finite unitary, without
constructing that certificate; general dynamic source rewind remains
open (Open problem~\ref{prob:tB-compiler}).
Set $E_0\psi=\psi\otimes m_0$ and $V=W^TE_0$.
For a specified unit vector $\eta\in\mathcal H_M$ the following are
equivalent:
\begin{enumerate}
\item There is a unitary $A$ on $Q$ such that $V\psi=A\psi\otimes\eta$
for every $\psi$.
\item $(\id_Q\otimes(\id_M-\proj\eta))V=0$.
\item The final $M$ marginal is $\proj\eta$ for every input density
operator on $Q$.
\end{enumerate}
Moreover the reduced channel $\rho\mapsto\tr_M(V\rho V^\dagger)$ is
unitary if and only if these conditions hold for some fixed pure $\eta$.
They continue to hold with an arbitrary untouched finite reference.
This disposal test is weaker than logical reset: $\eta$ need not equal
$m_0$. A clean implementation restores its workspace $Z$, but does not
by that fact restore the other logical interpreter factors.
Thus clock, position, journal, stack \emph{and all other retained
interpreter information} must collectively factor in a fixed pure state
if disposal is to implement a unitary on arbitrary $Q$.
Equal completion times alone are not a sufficient condition: the fixed
finite schedule of Theorem~\ref{thm:sA-four}, with keep target initially
zero, has one endpoint time while its final machine vectors
$\ket{00}$ and $\ket{10}$ remain correlated with the data, so collective
factorization fails there.
For a preserved computational-basis input with an orthonormal output
family, if
\[
 V\ket x=\ket{f_x}\ket{m_x},
\]
then an input matrix element $\ket x\bra y$ is sent after disposal to
$\langle m_y|m_x\rangle\ket{f_x}\bra{f_y}$. In particular orthogonal
clock/position records erase that off-diagonal element.
\end{theorem}
\begin{proof}
 \step{tB-clean-12}{Conditions 1 and 2 are equivalent.}
 \begin{proof}\pf\ Condition 1 implies 2 by projection. Under 2 set
 $A=(\id_Q\otimes\bra\eta)V$. Then $V\psi=A\psi\otimes\eta$
 for every $\psi$, and
 $A^\dagger A=V^\dagger V=\id$. Equal finite input and output
 dimensions make $A$ unitary.\end{proof}
 \step{tB-clean-23}{Conditions 2 and 3 are equivalent.}
 \begin{proof}\pf\ Condition 2 implies the stated pure marginal.
 Conversely apply 3 to every pure input. The squared norm of its
 component orthogonal to $Q\otimes\eta$ is the expectation of
 $\id\otimes(\id-\proj\eta)$, hence zero.\end{proof}
 \step{tB-clean-channel}{A unitary reduced channel implies a fixed pure
 environment vector.}
 \begin{proof}\pf\ If the reduced channel is conjugation by $A$,
 each pure output $V\ket i$ has $Q$ marginal
 $A\ket i\bra iA^\dagger$. The squared norm of its component
 orthogonal to $\operatorname{span}\{A\ket i\}\otimes\mathcal H_M$
 is the expectation of the complementary projection, hence zero.
 Therefore $V\ket i=A\ket i\otimes\eta_i$ for a unit vector
 $\eta_i$, on an orthonormal input basis. Density matrices span the
 full operator space: diagonal matrix units are pure projections and
 off-diagonal units are linear combinations of projections onto
 $(\ket i\pm\ket j)/\sqrt2$ and $(\ket i\pm i\ket j)/\sqrt2$.
 Channel equality thus holds on $\ket i\bra j$ and gives
 \[
 \langle\eta_j|\eta_i\rangle A\ket i\bra jA^\dagger
       =A\ket i\bra jA^\dagger.
 \]
 The rank-one operator on the right is nonzero, so the scalar is one.
 Hence $\|\eta_i-\eta_j\|^2
 =2-2\operatorname{Re}\langle\eta_j|\eta_i\rangle=0$.
 If $\dim Q=1$ there is only one vector. Conversely a common
 vector gives conjugation by $A$ on tracing.\end{proof}
 \step{tB-clean-elements}{The stated matrix-element formula holds.}
 \begin{proof}\pf\ Take the partial trace of
 $\ket{f_x}\bra{f_y}\otimes\ket{m_x}\bra{m_y}$.\end{proof}
 \qedstep
 \begin{proof}\pf\ Combine the steps. Tensoring the operator equality
 in condition 1 with the reference identity proves the reference claim.
 The equal-time failure follows from the nonzero, orthogonal machine
 terms and rank-two output proved in Theorem~\ref{thm:sA-four}.
 The condition is collective: constant individual mixed marginals do
 not establish factorization of the joint machine register.\end{proof}
\end{proof}

\begin{counterexample}[Exact undo at unequal ticks]\label{ex:tB-delay}
There are four qubits, in order $q,r,p_1,p_0$. The data are $Q=qr$;
$M=p_1p_0$ is a compact encoding of position, mode, the one-entry journal
and elapsed idle time. There is a single fixed mark, hence its identifier
and stack need zero additional bits. No separate clock or workspace is
used. Define $p=2p_1+p_0$, and the following complete permutations and
gate bits:
\[
\begin{array}{c|rrrr|rrrr}
 &\pi(0)&\pi(1)&\pi(2)&\pi(3)&a(0)&a(1)&a(2)&a(3)\\\hline
 q=0&1&2&3&0&1&1&0&0\\
 q=1&1&3&0&2&1&0&0&1
\end{array}
\]
The single fixed tick is
\[
 W\ket{q,r,p}=\ket{q,r\oplus a_q(p),\pi_q(p)}.          \tag{B5}
\]
At $p=0$ it performs a forward $X_r$ and records it in the position
encoding. For $q=0$ it undoes that entry at the second tick ($p=1$);
for $q=1$ it waits at $p=1$ and undoes at the third tick ($p=3$).
Completion is $p=2$ in either case. The transition $2\to3$ for $q=0$
is an idle tick recording elapsed completion time. Remaining transitions
complete the finite cycles; only $T=3$ ticks are claimed as execution.
For arbitrary amplitudes $\alpha,\beta$ and an arbitrary unit vector
$\chi\in\mathcal H_r$,
\[
 W^3\bigl((\alpha\ket0+\beta\ket1)_q\chi_r\ket0_p\bigr)
 =\alpha\ket0_q\chi_r\ket3_p+
   \beta\ket1_q\chi_r\ket2_p.                         \tag{B6}
\]
For $q=\ket+$ disposal of $p$ yields
$\tfrac12\id_q\otimes\proj\chi_r$, although each executed undo is
exact and $r$ is restored. The two completion ticks are 2 and 3.
\end{counterexample}
\begin{proof}
 \step{tB-delay-unit}{Equation (B5) is a permutation unitary.}
 \begin{proof}\pf\ For each unchanged $q$, $\pi_q$ is a permutation.
 From its output recover $p=\pi_q^{-1}(p')$ and then
 $r=r'\oplus a_q(p)$. This gives a unique inverse on all 16 basis states.
 Thus rule domains and ranges are disjoint.\end{proof}
 \step{tB-delay-evolve}{Equation (B6) holds.}
 \begin{proof}\pf\ For $q=0$, positions are $0,1,2,3$ and data gates
 are $X,X,\id$. For $q=1$, positions are $0,1,3,2$ and gates are
 $X,\id,X$. Both products are identity. Apply linearity. The inverse
 data action is $X^\dagger=X$. The compressed register $p$ is used
 only through the displayed total unitary (B5); no exhibition of this
 example as a certificate instance of Definition~\ref{def:tB-cert}
 is required for the disposal calculation.\end{proof}
 \qedstep
 \begin{proof}\pf\ The final machine states are orthogonal, so
 Theorem~\ref{thm:tB-clean} removes the cross terms. For $q=\ket+$
 the reduced $q$ state has purity $1/2$.\end{proof}
\end{proof}

\begin{lemma}[The global clock and absorbing completion]
Assume the clock initialization and tick (B3). Also let $U$ be any
unitary and $u,v$ any vectors in its Hilbert space.
For every initial state $\Phi\otimes\ket0_\Gamma$, the global clock
of (B3) is always the fixed basis
state $\ket T$ at the endpoint. It cannot itself cause the loss in (B6).
A data-dependent local elapsed-time register or position can do so.
Furthermore, if $Uu=v$ and $Uv=v$, then $u=v$.
\end{lemma}
\begin{proof}
 \step{tB-absorbing-clock}{At tick $t\leq T$, the clock factors as
 $\ket t$ from all other registers.}
 \begin{proof}\pf\ Induction in (B3) advances every input vector's
 clock by one, independent of its other registers. There is no wrap
 since $T<L_{\Gamma}$. A fixed pure factor has no correlations to discard.
 \end{proof}
 \step{tB-absorbing-injective}{$Uu=Uv$ implies $u=v$.}
 \begin{proof}\pf\ Multiply by $U^\dagger$.\end{proof}
 \qedstep
 \begin{proof}\pf\ Apply the preceding two steps; the local-position
 obstruction is Example~\ref{ex:tB-delay}.\end{proof}
\end{proof}

\begin{remark}[Checked halting and synchronization context]\label{rem:tB-halting}
Bernstein--Vazirani define exact halting by no final configurations before
$T$ and only final configurations at $T$
\cite[Definition 3.5.1, PDF p.~14]{BernsteinVazirani1997}.
Their synchronization theorem applies to polynomial-time classical
$f$ with output length depending only on input length. It supplies a
stationary normal-form reversible machine computing $(x,f(x))$ in time
depending only on input length
\cite[Theorem 4.1.3, PDF p.~17]{BernsteinVazirani1997}.
Its phrase ``whose running time depends only on the length of $x$''
does not supply a synchronization theorem for arbitrary coherent
mark selection after gates.

Linden--Popescu's specified auxiliary evolution retains elapsed stopping
time; they write ``ancilla contains a record of the time since the
computation halted''
\cite[PDF p.~2, equations (7)--(11) and text after (9)]{LindenPopescu1998}.
Example~\ref{ex:tB-delay} gives a finite bounded instance of the relevant
entanglement obstruction. We use neither a general nonexistence claim
nor a claim that unequal durations always destroy coherence.

Ozawa proves equality of monitored and unmonitored output probabilities
under his assumptions (I)--(III), including invariant halted tape output
\cite[PDF p.~2, equation (11); proof pp.~2--4,
equations (12)--(14) and following]{Ozawa1998}.
The phrase ``no more changes the halt qubit nor the tape string'' is
assumption (II), equation (6), on PDF p.~2. Equality of that output
distribution is distinct from unitary preservation of arbitrary quantum
output and does not contradict (B6).
Shi questions the composition argument for a particular universal-machine
proposal \cite[PDF pp.~7--8, Section 5, equations (12)--(14)]{Shi2002}.
This is contextual, not a theorem precluding bounded reversible circuits.
No halting paper carries a proof step of the results above.
\end{remark}

\subsection{Circuit costs and the unresolved compilation problem}

\begin{definition}[Explicit circuit cost parameters]\label{def:tB-cost}
For the following implementation bound, assume a supplied reversible
microcode circuit for (B3). It fetches a $w$-bit instruction from a
fixed classical table of length $\ell$ at an $a$-bit quantum address,
where $a=\lceil\log_2(\ell+1)\rceil$; out-of-range fetch returns a
specified default. This is a coherent circuit lookup, not an oracle
or assumed quantum random-access memory. It uses at most $\kappa$
journal accesses and at most $\kappa$ mark-stack accesses per tick,
where $\kappa$ is an absolute constant independent of
$\ell,a,w,N,d,u,v,\beta,g,b,h,z,q$ and $T$; it is part of the
constant hidden in $O(\cdot)$ in (B7). The cell width is at most $w$,
capacities $N$ and $d$, and binary cursor widths
$u=\lceil\log_2(N+1)\rceil$, $v=\lceil\log_2(d+1)\rceil$.
Take $w$ large enough for every descriptor and mark tuple.

Let $g$ be the sum of lengths of the supplied elementary action circuits
that are coherently selected during one tick, including all wire-address
selection, controlled gates, keeps and inverses. Let $\beta$ bound the
sum of sizes, including all input and output ports, of the other Boolean
circuits \emph{and their supplied inverse
routing computations}. These circuits must implement the certificate,
including behavior on invalid codes and completion. A table depending
on the clock is charged at its actual additional size in $\beta$ or $g$.
Let $h$ be all other logical machine bits and let $z$ bound reusable clean
workspace for the supplied component circuits. No small value of
$g,\beta,h,z$ follows from the existence of an abstract completion.
\end{definition}

\begin{theorem}[Conditional implementation cost]\label{thm:tB-cost}
Assume a certificate satisfying Definition~\ref{def:tB-cert}, its
specified tick and completion, and all supplied clean-circuit and
access hypotheses of Definition~\ref{def:tB-cost}. The following bound
is conditional on those data; it supplies no general dynamic source
compiler, which remains open (Open problem~\ref{prob:tB-compiler}).
A scan implementation of one tick has
\[
 G_{\rm tick}=O\bigl(\ell(a+w)+N(u+w)+d(v+w)
                  +\beta+g+b+w+1\bigr).                    \tag{B7}
\]
A sufficient total qubit count, including the input $q=|Q|$, is
\[
 q+h+(N+d)w+a+u+v+b+
 O(z+a+u+v+w+b+1).                                        \tag{B8}
\]
Here $h$ includes $A$, direction flags and retained routing records;
workspace for Boolean evaluation must already be included in $z$.
$T$ ticks cost $O(TG_{\rm tick})$ gates with the same wires.
The circuit always consists of exactly $T$ copies of the tick circuit.
A component in an idle rule sector still passes through that circuit;
no count of padded ticks for an individual sector component is subtracted
from the gate total. For a general input vector different sector
components may be idle at different ticks, so no single nonidle count
need exist. Separately, the static embedding of Theorem~\ref{thm:tB-A}
has $n$ scheduled logical-factor ticks and $T-n$ padding ticks, compared
with the direct Tier~A list of the $n$ logical factors $A_c$.
That schedule's factors are not counted as unit-cost elementary gates;
applying (B7) to any implementation requires all its circuit hypotheses.
A quantum-position interpreter therefore has
fetch and routing overhead even with classical control.
\end{theorem}
\begin{proof}
 \step{tB-cost-fetch}{Instruction lookup costs $O(\ell(a+w))$.}
 \begin{proof}\pf\ For each classical address $i$, reversibly compute
 the equality flag $[p=i]$ using $O(a)$ elementary Boolean gates and
 clean scratch, XOR the fixed $w$-bit word controlled by that flag,
 and undo the flag. Only one address can match. If no address matches,
 one further $O(a)$ range test and an $O(w)$ controlled XOR install
 the specified default word, after which the flag is uncomputed.
 Since $\ell\geq1$, this is absorbed by $O(\ell(a+w))$. Apply
 Lemma~\ref{lem:synthesis}; reversing the circuit erases a fetched
 descriptor when its inputs have been restored. A descriptor cannot
 be erased using a modified address without that restoration.\end{proof}
 \step{tB-cost-memory}{The stated journal and stack terms suffice.}
 \begin{proof}\pf\ Scan all $N$ journal cells, compute each cursor
 equality flag in $O(u)$ gates, and perform the required controlled
 bitwise XORs or swaps in $O(w)$ gates. Reverse the flag computation.
 The analogous scan of $d$ mark cells costs $O(d(v+w))$.
 The absolute access bound $\kappa$ absorbs repetitions of these scans.
 Cursor increments, decrements and comparisons use linear-size
 reversible binary arithmetic with retained carries uncomputed.
 The supplied microcode guarantees the required blank cells or retains
 their old values; a scan does not implement irreversible deletion.\end{proof}
 \step{tB-cost-rest}{The other terms and qubit count suffice.}
 \begin{proof}\pf\ Coherently control the supplied gate lists and
 evaluate the remaining Boolean circuits, charging $g$ and $\beta$.
 For the clock, evaluate the successor with a linear-size ripple
 circuit into a blank $b$-bit register, swap registers, and erase the
 old value by evaluating the predecessor into it. Boolean evaluation
 uses Lemma~\ref{lem:synthesis}; this costs $O(b+1)$ gates and clean
 workspace. Store the named cells and cursors, the $h$ other
 logical bits, a fetched word and a reusable set of equality/carry
 scratch bits. The clean component-workspace bound is $z$.
 This gives (B7)--(B8).\end{proof}
 \qedstep
 \begin{proof}\pf\ Repeat the same circuit $T$ times and reuse its
 zero-output workspace. For a superposition, there is
 one circuit consisting of $T$ ticks, not a sum of execution times over basis
 inputs.\end{proof}
\end{proof}

\begin{proposition}[Nesting costs in this semantics]\label{prop:tB-nesting}
Assume elementary signed journal entries, the certificate and explicit
tick-cost hypotheses of Theorem~\ref{thm:tB-cost}, and a fixed horizon
$T$. The bound contains no factor depending on nesting depth except
through its parameters $T,\ell,a,N,u,d,v,\beta,g,b,w,h,z$;
in particular it does not bound $T$ by a polynomial in the source
description. This conditional count supplies no general dynamic source
rewind construction, which remains open.
For the static family $P_0=a$ and
$P_{j+1}=\mathsf{mark}\{P_j\}\mathsf{rewind}\ K$, if the prescribed
expansion executes every elementary body gate and its adjoint and makes
no cancellations, its number $n_j$ of payload gate applications is
\[
 n_{j+1}=2n_j+k,\qquad n_j=2^jn_0+(2^j-1)k,                \tag{B9}
\]
where $n_0$ and $k$ are the lengths of the atom and keep circuits.
Any interpreter executing this list with at most one payload gate per
tick has $T\geq n_j$. This is a tight count for that expansion, not a
lower bound for all equivalent implementations. With $N=O(T)$ journal
capacity, (B7) gives the total scan bound
\[
 O\bigl(T^2(u+w)+
 T(\ell(a+w)+d(v+w)+\beta+g+b+w+1)\bigr),
\]
so no automatic linear-in-$T$ bound is obtained. Keeping a journal
does not itself remove (B9).
\end{proposition}
\begin{proof}
 \step{tB-nesting-global}{The total circuit has $T$ copies of its tick.}
 \begin{proof}\pf\ This is its definition and
 Theorem~\ref{thm:tB-cost}. Nesting affects the parameters and $T$,
 not an additional multiplication of this count.\end{proof}
 \step{tB-nesting-count}{Equation (B9) holds.}
 \begin{proof}\pf\ A marked expansion prints its body once, its
 $k$ keep gates once, and the adjoint body once. This gives the
 recurrence. Induction on $j$ gives the displayed solution.
 The circuit itself attains that count; serial execution cannot use
 fewer ticks than the number of gates it is required to execute.\end{proof}
 \qedstep
 \begin{proof}\pf\ Combine the steps and substitute $N=O(T)$ in
 (B7). No optimal-synthesis or general delayed-erasure theorem follows.
\end{proof}
\end{proof}

\subsection{Small exact instances and outstanding obligations}

\begin{example}[Finite matrices for independent checking]\label{ex:tB-table}
In each example below $p$ is a binary encoded cyclic instruction
counter, initially zero. It also encodes the fixed mode, journal
occupancy and fixed single-mark control; no unlisted auxiliary wire is
used. For a gate list $(G_0,\ldots,G_{L_{\Gamma}-1})$, notation
$[G_0,\ldots,G_{L_{\Gamma}-1}]$ means the complete unitary
$\sum_{p=0}^{L_{\Gamma}-1}\ket{p+1\bmod L_{\Gamma}}\bra p\otimes G_p$.
A fixed mark identifier needs no separate storage. These compressed
examples test algebraic and timing claims, not a variable-size stack
compiler.
\[
\begin{array}{c|c|c|c|l}
\text{case}&Q&M&T&\text{tick and endpoint from zero machine state}\\\hline
\mathrm{I}&q&p\ (1\text{ bit})&2&[X,X]:\ \psi\ket0\mapsto\psi\ket0\\
\mathrm{II}&q&p\ (2\text{ bits})&3&[H,X,H,\id]:\ \psi\ket0\mapsto Z\psi\ket3\\
\mathrm{III}&qr&p\ (2\text{ bits})&3&\text{(B5)--(B6); unequal completion ticks}\\
\mathrm{IV}&q&p\ (3\text{ bits})&7&[H,X,H,X,H,X,H,\id]:\
 \psi\ket0\mapsto -X\psi\ket7
\end{array}
\]
Case I uses two qubits, II three, and III--IV four. Case II is
$H^\dagger XH=Z$; case IV is an outer rewind with inner completed
operation $HXH=Z$ and outer keep $X$.
The minus sign in IV is part of the specified matrix and must not be
removed if the operation is later controlled.
\end{example}
\begin{proof}
 \step{tB-table-unit}{Each displayed tick is unitary.}
 \begin{proof}\pf\ Each counter input has a unique counter output,
 and its attached data gate is unitary. Case III was proved in
 Example~\ref{ex:tB-delay}.\end{proof}
 \step{tB-table-products}{The endpoints are as stated.}
 \begin{proof}\pf\ Apply the attached gates in counter order.
 $X^2=\id$, $HXH=Z$, and $ZXZ=-X$ by the two-by-two matrices.
 The counter ends at the indicated value.\end{proof}
 \qedstep
 \begin{proof}\pf\ The steps establish the table for every data
 input, including inputs entangled with an untouched reference.\end{proof}
\end{proof}

\begin{problem}[A compiler and a source simulation theorem]
\label{prob:tB-compiler}
For general bounded source programs with repeated dynamic mark visits,
construct the explicit coordinate encodings, local rules and efficient
completion in Definition~\ref{def:tB-cert}, and prove the reachable-domain
and disjoint-range obligations. This is not proved here. In particular,
after a keep changes a previously read input, newly executing the source
code can choose different instructions or fail to reach the saved visit.
It is not thereby the inverse of the recorded computation. A successful
translation must recover the exact old operation including its control,
or specify a different total completion and state the resulting semantic
change. Inspections of nonzero amplitudes, discarding a clock and taking
only a maximal deadline are not reversible Boolean operations on a
computational-basis state without additional records and a proof.
No equivalence with the executor's scheduling, rollback or budget rules
is established by this section.
\end{problem}

\begin{problem}[Synchronization, resources and unrestricted execution]
A source-level sufficient condition, proved for dynamic marks and arbitrary
intervening gates, which guarantees the factorization in
Theorem~\ref{thm:tB-clean}, remains open. Equal padded durations do not
clear the journal or mark stack. Extending the bounded construction to
an unknown unbounded runtime is not accomplished by the finite cyclic
clock. No efficient finite-alphabet synthesis of an arbitrary abstract
completion is asserted; exact synthesis of a prescribed arbitrary
completion need not exist over the finite alphabet, so a circuit claim
must supply its completion and circuit together, as in
Definition~\ref{def:tB-cert}. Finally, a general transformation
removing exponential repeated undo while retaining arbitrary quantum
keeps, with a proved time--space trade-off, is not supplied. The exact
count (B9) is not an impossibility result for such a transformation.
\end{problem}
